\documentclass[12pt]{amsart}

\usepackage{amsmath}
\usepackage{amssymb}
\usepackage{amsthm}
\usepackage{mathtools}
\usepackage{enumitem}
\usepackage[colorlinks=true,citecolor=blue,urlcolor=blue,linkcolor=blue]{hyperref}

% Theorem environments
\newtheorem{theorem}{Theorem}[section]
\newtheorem{lemma}[theorem]{Lemma}
\newtheorem{proposition}[theorem]{Proposition}
\newtheorem{corollary}[theorem]{Corollary}
\newtheorem{conjecture}[theorem]{Conjecture}

\theoremstyle{definition}
\newtheorem{definition}[theorem]{Definition}
\newtheorem{assumption}[theorem]{Assumption}
\newtheorem{example}[theorem]{Example}

\theoremstyle{remark}
\newtheorem{remark}[theorem]{Remark}

% Operators
\DeclareMathOperator{\tr}{tr}
\DeclareMathOperator{\diag}{diag}
\DeclareMathOperator{\rank}{rank}
\DeclareMathOperator{\MSE}{MSE}
\DeclareMathOperator{\relu}{ReLU}
\DeclareMathOperator{\spn}{span}
\DeclareMathOperator{\Var}{Var}
\DeclareMathOperator{\Fisher}{Fisher}
\DeclareMathOperator{\softmax}{softmax}

% Shortcuts
\newcommand{\R}{\mathbb{R}}
\newcommand{\E}{\mathbb{E}}
\newcommand{\Prob}{\mathbb{P}}
\newcommand{\bx}{\mathbf{x}}
\newcommand{\bX}{\mathbf{X}}
\newcommand{\bJ}{\mathbf{J}}
\newcommand{\bM}{\mathbf{M}}
\newcommand{\bU}{\mathbf{U}}
\newcommand{\bV}{\mathbf{V}}
\newcommand{\bI}{\mathbf{I}}
\newcommand{\bu}{\mathbf{u}}
\newcommand{\bv}{\mathbf{v}}
\newcommand{\bw}{\mathbf{w}}
\newcommand{\ba}{\mathbf{a}}
\newcommand{\bg}{\mathbf{g}}
\newcommand{\bp}{\mathbf{p}}
\newcommand{\bH}{\mathbf{H}}
\newcommand{\bF}{\mathbf{F}}
\newcommand{\bA}{\mathbf{A}}
\newcommand{\bB}{\mathbf{B}}
\newcommand{\by}{\mathbf{y}}
\newcommand{\bSigma}{\boldsymbol{\Sigma}}
\newcommand{\bGamma}{\boldsymbol{\Gamma}}
\newcommand{\bxi}{\boldsymbol{\xi}}
\newcommand{\bmu}{\boldsymbol{\mu}}
\newcommand{\norm}[1]{\left\lVert #1 \right\rVert}
\newcommand{\abs}[1]{\left\lvert #1 \right\rvert}
\newcommand{\ip}[2]{\langle #1, #2 \rangle}

\title[Feature Learning, Effective Dimension, and Critical Depth]{%
  Feature Learning, Effective Dimension, and\\
  Architecture-Specific Critical Depth}

\author{Lightman Chang}
\address{Independent Researcher}
\email{lightman.chang@gmail.com}

\date{\today}

\subjclass[2020]{Primary 68T07; Secondary 62G05, 41A25, 60J60}

\keywords{Feature learning, neural tangent kernel, effective dimension,
  curse of dimensionality, single-index models, critical depth,
  ResNet, Transformer}

\begin{document}

\begin{abstract}
The minimax rate for nonparametric regression of $C^{s}$ functions on
$\R^{d}$ is $n^{-2s/(2s+d)}$, which deteriorates rapidly with the ambient
dimension. Practical neural networks trained on high-dimensional inputs
nonetheless generalize at rates that do not exhibit such severe dimension
dependence. We give a precise mechanistic account of this gap for the
\emph{single-index} target class $f^{*}(\bx) = \rho(\langle \bv^{*}, \bx \rangle)$
under Gaussian inputs, working within the noise-propagation framework
of~\cite{paperA}. Our principal contribution (Theorem~\ref{thm:fl-collapse})
shows that, conditional on the alignment hypothesis established
in~\cite{paperB}, the test risk of two-layer ReLU networks transitions from
the NTK-regime rate $\Theta(n^{-2s/(2s+d)})$ to
$\Theta(n^{-2s/(2s+1)})$ in the feature-learning regime; the proof routes
the NTK rate through spherical-harmonic multiplicities $N_{l} \asymp l^{d-2}$
and the feature-learning rate through the rank-one effective dimension
induced by the collapsed equivalent kernel. We then introduce an
architecture-specific critical depth $L^{*}$ (Theorem~\ref{thm:critical-depth})
defined as the smallest depth at which the effective rank of the noise
propagation operator falls below $n^{2s/(2s+1)}$, and we derive
$L^{*}_{\mathrm{FC}} = \lceil(d-1)/2\rceil$ for fully-connected ReLU
networks (with the per-layer exponent $\alpha_{0} = (d+1)/(2(d-1))$ from
spherical-harmonic NTK eigenvalues, tracking depth-dependent constants
explicitly) and a depth-decoupled formula for residual networks. The
analogous result for transformer architectures is stated as a heuristic
conjecture (Conjecture~\ref{conj:transformer}) pending a rigorous
attention-head spectral computation. The cross-entropy extension via a
Fisher-corrected operator $\bM_{\mathrm{CE}}$ and the rank-$(K-1)$ effective
rank at the Neural Collapse fixed point are treated in the companion
paper~\cite{paperD}.
\end{abstract}

\maketitle

%% ============================================================
\section{Introduction}\label{sec:intro}
%% ============================================================

\subsection{The dimension dependence of the kernel rate}

A long-standing tension in the theory of overparameterized learning concerns
the role of the ambient input dimension $d$. Classical nonparametric
estimation gives a sharp lower bound: for the H\"older class
$C^{s}(\R^{d})$, no estimator can attain a rate better than
$n^{-2s/(2s+d)}$, and kernel ridge regression with a smooth rotation-invariant
kernel attains exactly this rate. The exponent is dominated by $d$ for any
fixed smoothness $s$, so that even modest inputs (e.g.\ $d=100$) require an
exponentially large number of samples for moderate accuracy. This is the
\emph{curse of dimensionality}.

Empirical practice presents a different picture. Networks trained on inputs
of dimension $d$ in the hundreds or thousands typically attain test errors
that scale far better than $n^{-2s/(2s+d)}$. Two complementary explanations
are usually offered. The first is structural: real targets are well
approximated by low-complexity composites (single-index, multi-index, or
sparse functions of low intrinsic dimension), so that the relevant rate is
governed by an intrinsic dimension $k \ll d$. The second is algorithmic:
gradient-based training does not produce a kernel-regime estimator; rather, it
performs feature learning, and the post-training Jacobian concentrates
spectral mass on directions that are aligned with the target structure.

The first explanation is essentially structural and rests on assumptions about
the data-generating process. The second, however, makes a more delicate
prediction: even when the input is genuinely $d$-dimensional, an
overparameterized network can attain a rate that scales as if the input were
one-dimensional, provided that feature learning succeeds at aligning the
weights with the target direction. The present paper makes this prediction
quantitative and unconditional, given the alignment guarantees already
established in~\cite{paperB}.

\subsection{Setting and prior work}

We work within the framework introduced in~\cite{paperA}, which models the
test mean squared error of an interpolating predictor through the
\emph{noise propagation operator} $\bM = \bSigma^{-1} \bGamma\, \bSigma^{-1}$,
where $\bSigma$ is the diagonal matrix of singular values of the training
Jacobian and $\bGamma$ is the test--train cross-correlation matrix of
gradient features. The trace $\tr(\bM)$ governs the noise contribution to
test error, and the benign overfitting trichotomy
($\beta < \alpha + 1/2$ catastrophic, $\beta = \alpha + 1/2$ tempered,
$\beta > \alpha + 1/2$ benign) follows from the power-law decay of singular
values $\sigma_{j} \asymp j^{-\alpha}$ and gradient-feature norms
$\varrho_{j} \asymp j^{-\beta}$. The companion paper~\cite{paperB} establishes
the alignment hypothesis: in the rich regime, gradient flow drives the
weights of every active neuron of a two-layer ReLU network into the span of
the teacher direction $\bv^{*}$.

We refer to~\cite{paperA} for the unified generalization bound and the
benign/tempered/catastrophic dichotomy, and to~\cite{paperB} for the
alignment lemma and its proof via ReLU directional independence. The
present paper does \emph{not} reprove these results: we cite them as
black-box inputs and develop their consequences.

The literature on dimension-free rates for neural networks is now
substantial. Bach~\cite{Bach2017} gave a convex/variational analysis of two-layer
networks and derived dimension-adaptive approximation rates. Bietti and
Bruna~\cite{BiettiBruna2021} studied spherical harmonic decompositions of
the NTK and gave matching upper and lower bounds on kernel ridge regression.
Damian, Lee, and Soltanolkotabi~\cite{DLS2022} proved that gradient
descent learns single-index targets via a feature-learning step that
escapes the kernel regime, attaining the dimension-free rate
$n^{-2s/(2s+1)}$.
Ben Arous, Gheissari, and Jagannath~\cite{BenArous2021} characterized
the sample complexity of online SGD for single-index recovery in terms
of the information exponent $\kappa$, identifying the $\Theta(d^{\kappa})$
sample threshold for signal detection. The two results are complementary:
the latter analyzes the signal-detection phase (number of samples needed
to escape from the kernel regime), while the former analyzes the rate
once detection has occurred.
Rakhlin and Zhai~\cite{RZ2019} showed that interpolation by Laplace kernels
is consistent only in high dimension, isolating the dimension dependence of
benign overfitting in the NTK regime. Papyan, Han, and
Donoho~\cite{papyan2020} introduced the Neural Collapse phenomenon, in
which last-layer features and class means collapse onto a simplex
equiangular tight frame at the terminal phase of training; the
cross-entropy implications of Neural Collapse for our framework are
treated in the companion paper~\cite{paperD}.

\subsection{Contributions}

This paper makes two main contributions, both conditional on the
alignment hypothesis of~\cite{paperB}, plus a heuristic conjecture for
transformers.

\begin{enumerate}[label=(\roman*)]
\item \textbf{Feature-learning dimension collapse}
  (Theorem~\ref{thm:fl-collapse}). For single-index targets with smoothness
  $s$, we show that the test risk transitions from the NTK-regime rate
  $\Theta(n^{-2s/(2s+d)})$ to the dimension-free rate
  $\Theta(n^{-2s/(2s+1)})$ when feature learning succeeds. The proof
  identifies the spherical-harmonic multiplicity $N_{l} \asymp l^{d-2}$ as
  the source of the curse of dimensionality in the NTK regime, and
  the multiplicity collapse to a one-dimensional kernel (at the level of
  the population equivalent kernel) as the mechanism responsible for the
  dimension-free rate after alignment.

\item \textbf{Architecture-specific critical depth for FC and ResNet}
  (Theorem~\ref{thm:critical-depth}). We define $L^{*}_{A}$ as the smallest
  depth at which the effective rank $r_{\mathrm{eff}}(\bM_{A})$ falls
  below $n^{2s/(2s+1)}$. Using the per-layer singular-value exponent
  $\alpha_{0} = (d+1)/(2(d-1))$ from Lemma~\ref{lem:alpha0}, we derive
  $L^{*}_{\mathrm{FC}} = \lceil(d-1)/2\rceil$ for fully-connected ReLU
  stacks (via depth-multiplicative spectral decay; the depth-dependent
  constant $L^{L\alpha_{0}}$ is tracked explicitly in
  Proposition~\ref{prop:depth-mult}), and
  $L^{*}_{\mathrm{ResNet}} \in \{1,\infty\}$ decoupled from depth (via the
  identity-bypass structure of skip connections).
  For transformers, we record the analogous heuristic prediction as an
  explicit conjecture (Conjecture~\ref{conj:transformer}) and discuss
  what would be required to make it rigorous.
\end{enumerate}

The cross-entropy extension via a Fisher-corrected operator
$\bM_{\mathrm{CE}}$ and the rank-$(K-1)$ effective rank at the Neural
Collapse fixed point are treated separately in the companion
paper~\cite{paperD}.

\subsection{Delta from prior work}\label{sec:delta}

The kernel-regime rate $n^{-2s/(2s+d)}$ for kernel ridge regression in
the NTK regime is folklore (see Caponnetto--De Vito~\cite{capDeVito2007}
for the abstract minimax-rate theory and the references therein for
spherical NTK derivations). The dimension-free single-index rate
$n^{-2s/(2s+1)}$ is established for one-pass SGD by Damian--Lee
--Soltanolkotabi~\cite{DLS2022} and for one-step gradient by
Ba--Mei--Misiakiewicz--Montanari~\cite{ba2022mei}. The new contribution
of this paper is the systematic identification of the spherical-harmonic
multiplicity $N(d,l) \asymp l^{d-2}$ as the precise source of the curse
in the kernel regime, the rank-one effective dimension as the mechanism
of the rate under alignment, and the depth-multiplicative chain that gives
$L^{*}_{\mathrm{FC}} = \Theta(d)$ for plain ReLU stacks. The
ResNet decoupling and the transformer conjecture are, to our knowledge,
not in the prior literature in this exact form.

The multi-index extension lies adjacent to the staircase-learning
literature of Abbe--Adsera--Misiakiewicz~\cite{abbe2022} and
Mei--Misiakiewicz--Montanari~\cite{MMM2022}; we discuss this in
Section~\ref{sec:discussion}.

\subsection{Organization}

Section~\ref{sec:prelim} recalls the noise propagation framework
of~\cite{paperA} and the alignment hypothesis of~\cite{paperB}.
Section~\ref{sec:fl-statement} states Theorem~\ref{thm:fl-collapse}, the
feature-learning dimension collapse. Sections~\ref{sec:fl-ntk}
and~\ref{sec:fl-fl} prove the NTK-regime and feature-learning halves of
Theorem~\ref{thm:fl-collapse} respectively. Section~\ref{sec:ld-statement}
states Theorem~\ref{thm:critical-depth}; Sections~\ref{sec:ld-fc} and
\ref{sec:ld-resnet} prove the FC and ResNet cases.
Section~\ref{sec:ld-transformer} states Conjecture~\ref{conj:transformer}
for transformers. Section~\ref{sec:discussion} discusses limitations and
open directions.

%% ============================================================
\section{Preliminaries}\label{sec:prelim}
%% ============================================================

\subsection{Recap of the noise propagation framework}

We follow the notation of~\cite{paperA}. Let $f(\cdot\,;\theta) \colon \R^{d} \to \R$
be a parametric predictor with parameters $\theta \in \R^{p}$, and let
$S = \{(\bx_{i}, y_{i})\}_{i=1}^{n}$ be a training set with
$y_{i} = g^{*}(\bx_{i}) + \xi_{i}$, $\xi_{i}$ i.i.d.\ sub-Gaussian with
variance proxy $\sigma^{2}$. Let $\theta^{*}$ be a minimum-norm
interpolating parameter and let
\[
  \bJ \in \R^{n \times p}, \qquad
  \bJ_{i,:} = \nabla_{\theta} f(\bx_{i};\theta^{*})^{\top},
\]
be the training Jacobian, with SVD $\bJ = \bU \bSigma \bV^{\top}$. The
$j$-th gradient feature is $\psi_{j}(\bx) = \nabla_{\theta} f(\bx;\theta^{*})^{\top}\bv_{j}$,
the cross-correlation matrix is
$\bGamma_{jl} = \E_{\bx}[\psi_{j}(\bx)\psi_{l}(\bx)]$, and the
noise propagation operator is
\[
  \bM = \bSigma^{-1} \bGamma\, \bSigma^{-1}
  \in \R^{n \times n}.
\]
Theorem~1 of~\cite{paperA} states that for any $\delta \in (0,1)$, with
probability $1-\delta$,
\begin{equation}\label{eq:gen-bound}
  \MSE_{\mathrm{test}}
  \leq C \bigl( B_{\mathrm{signal}}^{2}
       + \sigma^{2} \tr(\bM)
       + \sigma^{2} \norm{\bM}_{F} \sqrt{\log(1/\delta)}
       + \delta_{\mathrm{lin}}^{2} \bigr).
\end{equation}
Theorem~2 of~\cite{paperA} states that under power-law decay
$\sigma_{j} \asymp j^{-\alpha}$, $\varrho_{j} \asymp j^{-\beta}$, with
$\varrho_{j} := \sqrt{\bGamma_{jj}}$ (the diagonal feature-norm sequence),
\[
  \tr(\bM) < \infty \iff \beta > \alpha + \tfrac{1}{2}.
\]

\begin{definition}[Effective rank]\label{def:reff}
For a positive semidefinite operator $\bA$ with non-negative eigenvalues
$\lambda_{1} \geq \lambda_{2} \geq \cdots$, the \emph{effective rank} is
\[
  r_{\mathrm{eff}}(\bA) := \frac{(\sum_{j} \lambda_{j})^{2}}{\sum_{j} \lambda_{j}^{2}}
  = \frac{(\tr \bA)^{2}}{\norm{\bA}_{F}^{2}}.
\]
For $\bM$, the effective rank quantifies the number of significant
noise-propagation directions; benign overfitting requires $r_{\mathrm{eff}}(\bM)$
to be at most polylogarithmic in $n$.
\end{definition}

\subsection{Single-index targets and the alignment hypothesis}

\begin{definition}[Single-index target]\label{def:single-index}
A function $g^{*} \colon \R^{d} \to \R$ is a \emph{single-index target}
with link $\rho$ and direction $\bv^{*}$ if there exist
$\bv^{*} \in \mathbb{S}^{d-1}$ and $\rho \colon \R \to \R$ such that
\[
  g^{*}(\bx) = \rho(\langle \bv^{*}, \bx \rangle)
  \qquad \text{for all } \bx \in \R^{d}.
\]
We say $g^{*} \in \mathcal{H}^{s}$ if $\rho \in C^{s}(\R)$ with bounded
derivatives up to order $s$.
\end{definition}

We work throughout under the standard Gaussian input distribution
$\bx \sim \mathcal{N}(0, \bI_{d})$, with link function $\rho$ that is
non-affine and has information exponent $\kappa = 1$ in the sense
of~\cite{BenArous2021}; that is, the first non-trivial Hermite coefficient
of $\rho$ is non-zero. This is the regime in which the alignment
hypothesis of~\cite{paperB} applies.

\begin{assumption}[Alignment, after \cite{paperB}]\label{ass:alignment}
Let $f_{\theta}(\bx) = p^{-1/2} \sum_{j=1}^{p} a_{j}\, \relu(\bw_{j}^{\top}\bx)$
be a two-layer ReLU student. Suppose the inputs are
$\bx \sim \mathcal{N}(0,\bI_{d})$ and the target is single-index
$g^{*}(\bx) = \rho(\langle \bv^{*}, \bx\rangle)$ with $\kappa = 1$. Suppose
$f_{\theta}$ is trained by population gradient flow from a small-scale
initialization. Then at any limit point $\theta^{*}$ of the gradient flow,
the active neurons satisfy $\bw_{j} \in \spn(\bv^{*})$, and the dead
neurons satisfy $\norm{\bw_{j}} = O(\alpha_{\mathrm{init}})$.
\end{assumption}

This is exactly the conclusion of Theorem~5a--c of~\cite{paperB}; we use it
as a black box. The point of Theorem~\ref{thm:fl-collapse} below is to
quantify the consequence of Assumption~\ref{ass:alignment} for the test
risk.

\subsection{Equivalent kernels in the NTK and feature-learning regimes}

\begin{definition}[Equivalent kernel]\label{def:eq-kernel}
The \emph{equivalent kernel} of an interpolating predictor at $\theta^{*}$
is the population kernel
\[
  K_{\theta^{*}}(\bx, \bx') := \E_{\theta^{*}}\bigl[
    \nabla_{\theta} f(\bx; \theta^{*})^{\top}
    \nabla_{\theta} f(\bx'; \theta^{*}) \bigr],
\]
where the expectation is taken over the randomness in the trained
parameters (e.g.\ over initializations or training noise). When
$f_{\theta}$ is in the NTK regime, $K_{\theta^{*}} \approx K_{\mathrm{NTK}}$
is the rotation-invariant Neural Tangent Kernel of \cite{JGH2018}. After
feature learning, $K_{\theta^{*}}$ is no longer rotation-invariant.
\end{definition}

\begin{lemma}[Spherical harmonic decomposition of the NTK; \cite{BiettiBruna2021}]
\label{lem:ntk-spec}
Let $K_{\mathrm{NTK}}$ be the NTK of a two-layer ReLU network with input
dimension $d$ and Gaussian inputs. Then $K_{\mathrm{NTK}}$ admits the
spherical-harmonic eigendecomposition
\[
  K_{\mathrm{NTK}}(\bx, \bx')
  = \sum_{l=0}^{\infty} \lambda_{l}^{(d)}
    \sum_{m=1}^{N(d,l)} Y_{l,m}(\bx) Y_{l,m}(\bx'),
\]
where $\{Y_{l,m}\}$ is the orthonormal basis of spherical harmonics of
degree $l$ on $\mathbb{S}^{d-1}$, $N(d,l) = \binom{d+l-1}{d-1} - \binom{d+l-3}{d-1}$
is the multiplicity, and there exist constants $c_{l}, C_{l} > 0$ such
that, for $l \geq 1$,
\[
  c_{l}\, l^{-(d+1)} \leq \lambda_{l}^{(d)} \leq C_{l}\, l^{-(d+1)}.
\]
The asymptotic multiplicity is $N(d,l) = \Theta(l^{d-2})$ as $l \to \infty$
with $d$ fixed.
\end{lemma}

\begin{remark}\label{rem:cite-eq-kernel}
We will use Lemma~\ref{lem:ntk-spec} as a black box; the eigenvalue rate
$l^{-(d+1)}$ and the multiplicity $l^{d-2}$ are standard consequences of
the homogeneity properties of the ReLU activation and the structure of the
spherical Funk--Hecke formula. See~\cite{BiettiBruna2021}, Theorem~3, for
a complete proof.
\end{remark}

%% ============================================================
\section{Theorem 1: Feature Learning Dimension Collapse}\label{sec:fl-statement}
%% ============================================================

\begin{theorem}[Feature Learning Dimension Collapse]\label{thm:fl-collapse}
Let $g^{*}(\bx) = \rho(\langle \bv^{*}, \bx \rangle)$ be a single-index target
with $\rho \in \mathcal{H}^{s}$ for some $s > 1/2$, $\bv^{*} \in \mathbb{S}^{d-1}$,
and let $\bx \sim \mathcal{N}(0, \bI_{d})$. Let $f_{\theta}$ be a two-layer
ReLU network of width $p$ trained on $S \sim (g^{*} + \mathcal{N}(0,\sigma^{2}))^{n}$
to a minimum-norm interpolant $\theta^{*}$.

\textbf{(NTK regime).} If the network is trained in the NTK regime so that
the equivalent kernel coincides with the population NTK in the sense that
$\bM \approx \bM_{\mathrm{NTK}}$ to leading order, the optimally
ridge-regularized test risk satisfies
\begin{equation}\label{eq:rate-ntk}
  \E_{S} \MSE_{\mathrm{test}}^{\mathrm{NTK}}(n)
  = \Theta(n^{-2s/(2s+d)}).
\end{equation}

\textbf{(Feature-learning regime).} If the network is trained in the rich
regime so that Assumption~\ref{ass:alignment} holds at $\theta^{*}$, the
optimally ridge-regularized test risk satisfies
\begin{equation}\label{eq:rate-fl}
  \E_{S} \MSE_{\mathrm{test}}^{\mathrm{FL}}(n)
  = \Theta(n^{-2s/(2s+1)}).
\end{equation}

The improvement factor is
\[
  \frac{\MSE^{\mathrm{NTK}}}{\MSE^{\mathrm{FL}}}
  = \Theta\!\bigl(n^{2s(d-1)/((2s+1)(2s+d))}\bigr),
\]
which grows polynomially in $n$ for any fixed $d > 1$.
\end{theorem}

\begin{remark}[Status of the alignment hypothesis]\label{rem:status}
The condition under which the feature-learning conclusion applies is
exactly that of Assumption~\ref{ass:alignment}, which is in turn
established by Theorem~5a--c of~\cite{paperB}. The present theorem is thus
a conditional statement of the form ``alignment $\Rightarrow$ dimension-free
rate''. Removing the alignment hypothesis (e.g.\ proving the same rate
without an alignment guarantee) is open; see Section~\ref{sec:discussion}.
\end{remark}

The proof of Theorem~\ref{thm:fl-collapse} occupies the next two sections.
Section~\ref{sec:fl-ntk} proves the NTK rate~\eqref{eq:rate-ntk};
Section~\ref{sec:fl-fl} proves the feature-learning rate~\eqref{eq:rate-fl}.

%% ============================================================
\section{Proof of Theorem~\ref{thm:fl-collapse}: NTK rate via spherical
  harmonic multiplicity}\label{sec:fl-ntk}
%% ============================================================

We prove~\eqref{eq:rate-ntk} by reducing the test risk to that of optimally
ridge-regularized regression with respect to the population NTK and
applying classical kernel-rate arguments.

\subsection{Reduction to kernel ridge regression}

In the NTK regime, the linearization residual $\delta_{\mathrm{lin}}$
in~\eqref{eq:gen-bound} is $o(1)$ and the gradient features
$\psi_{j}(\bx) = \nabla_{\theta} f(\bx;\theta^{*})^{\top}\bv_{j}$ depend
only on the kernel structure. Specifically, the columns of $\bV$
diagonalize $\bJ^{\top}\bJ$, which converges in the wide limit to the
Gram matrix of $K_{\mathrm{NTK}}$ on the training set. The unified bound
\eqref{eq:gen-bound} therefore reduces, up to $o(1)$ corrections, to the
test risk of kernel ridge regression with kernel $K_{\mathrm{NTK}}$.

\subsection{Source-condition formulation}

Let $\{(\lambda_{l,m}, e_{l,m})\}$ enumerate the eigenpairs of
$K_{\mathrm{NTK}}$ in the spherical-harmonic basis of
Lemma~\ref{lem:ntk-spec}. The target
$g^{*}(\bx) = \rho(\langle \bv^{*}, \bx \rangle)$, viewed as a function on
$\mathbb{R}^{d}$ with Gaussian measure, has Hermite expansion
\[
  \rho(\langle \bv^{*}, \bx \rangle)
  = \sum_{l=0}^{\infty} c_{l}\, H_{l}(\langle \bv^{*}, \bx \rangle),
\]
where $H_{l}$ is the $l$-th (probabilist's) Hermite polynomial. For a
$C^{s}$ link function $\rho$, the standard Hermite-decay result is
$|c_{l}|^{2} = O(l^{-2s-1/2})$ (where the additional $1/2$ relative to
the Fourier case $|c_{l}|^{2} = O(l^{-2s})$ reflects the
$(2^{l}\,l!)^{-1/2}$ Hermite normalization on Gaussian measure; see e.g.\
\cite[Theorem 4.6]{AtkinsonHan2012}). In particular,
$\sum_{l} c_{l}^{2}\, l^{2s-1/2} < \infty$ for $\rho \in C^{s}(\R)$,
which is the source condition we use in the bias bound below.

Each Hermite mode $H_{l}(\langle \bv^{*}, \bx\rangle)$ is supported on the
spherical-harmonic shell of degree $l$ in the eigenbasis of $K_{\mathrm{NTK}}$;
indeed, $H_{l}(\langle \bv^{*}, \bx\rangle)$ is a degree-$l$ polynomial
that depends only on $\langle\bv^{*},\bx\rangle$, hence is a degree-$l$
zonal harmonic with respect to $\bv^{*}$. By rotational invariance of
$K_{\mathrm{NTK}}$, this zonal harmonic is a single eigenfunction of
$K_{\mathrm{NTK}}$ with eigenvalue $\lambda_{l}^{(d)}$. Consequently,
in the eigenbasis of $K_{\mathrm{NTK}}$, the coefficients of $g^{*}$ are
supported on $N(d,l) \asymp l^{d-2}$ many shells of equal magnitude up to
the zonal/non-zonal split.

\subsection{Bias--variance tradeoff and the optimal rate}

For ridge regression with regularization $\eta > 0$ and i.i.d.\ Gaussian
inputs and noise, the standard bias--variance bound (see, e.g.,
\cite[Section~2]{RZ2019}) gives
\begin{equation}\label{eq:bv}
  \E_{S} \MSE_{\mathrm{test}}
  \lesssim \underbrace{\sum_{l,m} \frac{\eta^{2} c_{l,m}^{2}}{(\lambda_{l}^{(d)} + \eta)^{2}}}_{\text{bias}}
  + \underbrace{\frac{\sigma^{2}}{n} \sum_{l,m} \frac{\lambda_{l}^{(d)}}{\lambda_{l}^{(d)} + \eta}}_{\text{variance}},
\end{equation}
where $c_{l,m}$ are the eigen-coefficients of $g^{*}$.

\textbf{Bias term.}
By the source condition $\sum_{l} c_{l}^{2}\, l^{2s} < \infty$ and the
upper bound $\lambda_{l}^{(d)} \leq C l^{-(d+1)}$ from
Lemma~\ref{lem:ntk-spec},
\[
  \sum_{l,m} \frac{\eta^{2} c_{l,m}^{2}}{(\lambda_{l}^{(d)} + \eta)^{2}}
  \asymp \eta^{2s/(d+1)},
\]
where the equivalence follows by partitioning the sum at the index
$l_{\eta}$ for which $\lambda_{l_{\eta}}^{(d)} = \eta$, namely
$l_{\eta} \asymp \eta^{-1/(d+1)}$, and integrating both contributions.

\textbf{Variance term.}
The effective dimension is
\[
  \mathrm{df}(\eta) := \sum_{l,m} \frac{\lambda_{l}^{(d)}}{\lambda_{l}^{(d)}+\eta}
  \asymp \sum_{l \leq l_{\eta}} N(d,l)
  \asymp l_{\eta}^{d-1}
  \asymp \eta^{-(d-1)/(d+1)},
\]
where we used $N(d,l) \asymp l^{d-2}$ from Lemma~\ref{lem:ntk-spec} and
summed via $\sum_{l \leq L} l^{d-2} \asymp L^{d-1}$. The variance term is
therefore $\asymp \sigma^{2}\, \eta^{-(d-1)/(d+1)}/n$.

\textbf{Combining.}
Setting bias and variance equal,
$\eta^{2s/(d+1)} \asymp \sigma^{2}\, \eta^{-(d-1)/(d+1)}/n$,
yields $\eta^{*} \asymp n^{-(d+1)/(2s+d)}$ and
\[
  \E_{S} \MSE_{\mathrm{test}}^{\mathrm{NTK}}(n)
  \asymp (\eta^{*})^{2s/(d+1)}
  \asymp n^{-2s/(2s+d)}.
\]
The matching lower bound follows from the standard nonparametric minimax
lower bound on $\mathcal{H}^{s}(\mathbb{S}^{d-1})$, since the NTK at this
optimally tuned ridge attains the minimax rate; see~\cite{RZ2019}. This
proves~\eqref{eq:rate-ntk}.\qed

\subsection{Where the dimension dependence enters}

The exponent $d$ in $n^{-2s/(2s+d)}$ comes from two compounding sources in
the variance term: the eigenvalue rate $l^{-(d+1)}$ and the multiplicity
$N(d,l) = \Theta(l^{d-2})$. The product $\lambda_{l}^{(d)} N(d,l) \asymp l^{-3}$
(after summation) is dimension-independent, but the threshold $l_{\eta}$
scales as $\eta^{-1/(d+1)}$ and is summed against $N(d,l)$, producing
$\mathrm{df}(\eta) \asymp \eta^{-(d-1)/(d+1)}$. The exponent $(d-1)/(d+1)$
is the source of the curse: when $d$ is large, the effective dimension
grows almost linearly with $\eta^{-1}$, and the variance dominates.

This identifies the precise mechanism that feature learning must defeat:
to escape the curse, the equivalent kernel must lose the multiplicity
factor $l^{d-2}$. The next section shows how alignment achieves this.

%% ============================================================
\section{Proof of Theorem~\ref{thm:fl-collapse}: Feature-learning rate via
  rank-one effective dimension}\label{sec:fl-fl}
%% ============================================================

We now prove~\eqref{eq:rate-fl}. The argument proceeds in three steps:
(i) characterize the equivalent kernel under alignment as a one-dimensional
ridge kernel along $\bv^{*}$; (ii) compute the resulting eigenvalue
sequence and effective dimension; (iii) optimize the bias--variance
tradeoff.

\subsection{Equivalent kernel under alignment}

\begin{lemma}[Population kernel collapse]\label{lem:kernel-collapse}
Suppose Assumption~\ref{ass:alignment} holds at $\theta^{*}$, and let
$\bar K(t,t') := K_{\mathrm{NTK}}^{(1)}(t,t')$ denote the one-dimensional
ReLU NTK on $\R$. Define the \emph{population equivalent kernel}
$\bar K^{\mathrm{pop}}_{\theta^{*}}(\bx,\bx')$ by integrating out the
component orthogonal to $\bv^{*}$:
\[
  \bar K^{\mathrm{pop}}_{\theta^{*}}(\bx,\bx')
  := \E_{\bx_{\perp},\bx'_{\perp}}\!\bigl[\,K_{\theta^{*}}(\bx,\bx')\bigr],
\]
where $\bx = \langle\bv^{*},\bx\rangle\,\bv^{*} + \bx_{\perp}$ and likewise
for $\bx'$, with $\bx_{\perp},\bx'_{\perp}$ drawn from the Gaussian marginal
on the orthogonal complement of $\bv^{*}$. Then
\[
  \bar K^{\mathrm{pop}}_{\theta^{*}}(\bx, \bx')
  = \bar K\bigl(\langle \bv^{*}, \bx\rangle, \langle \bv^{*}, \bx'\rangle\bigr)
  + c_{0} + O(\alpha_{\mathrm{init}}^{2}),
\]
where $c_{0}$ is a constant that does not depend on $\bx,\bx'$ (and is
absorbed into the constant ``bias'' direction of the equivalent kernel).
\end{lemma}

\begin{proof}
Under Assumption~\ref{ass:alignment}, the active neurons satisfy
$\bw_{j} = w_{j}\bv^{*}$ for some scalar $w_{j} \in \R$, and dead neurons
have norm $O(\alpha_{\mathrm{init}})$. The gradient of the predictor with
respect to its parameters at $\bx$ is therefore
\begin{align*}
  \nabla_{a_{j}} f_{\theta^{*}}(\bx) &= p^{-1/2} \relu(w_{j} \langle\bv^{*},\bx\rangle), \\
  \nabla_{\bw_{j}} f_{\theta^{*}}(\bx) &= p^{-1/2} a_{j} \mathbf{1}\{w_{j}\langle\bv^{*},\bx\rangle > 0\} \bx,
\end{align*}
for active $j$, with both quantities $O(\alpha_{\mathrm{init}})$ for dead $j$.
The pointwise kernel $K_{\theta^{*}}(\bx,\bx')$ contains a contribution from
the $a_{j}$-gradients that is a function of
$(\langle\bv^{*},\bx\rangle,\langle\bv^{*},\bx'\rangle)$ only, plus a
contribution from the $\bw_{j}$-gradients that is the same function multiplied
by $\bx^{\top}\bx'$. Decomposing
$\bx^{\top}\bx' = \langle\bv^{*},\bx\rangle\langle\bv^{*},\bx'\rangle + \bx_{\perp}^{\top}\bx'_{\perp}$,
the first term is again a function of $(\langle\bv^{*},\bx\rangle,\langle\bv^{*},\bx'\rangle)$.
The second term, $\bx_{\perp}^{\top}\bx'_{\perp}$, is a random quantity that
depends on the perpendicular components.

Integrating out $\bx_{\perp}$ and $\bx'_{\perp}$ (which are independent
mean-zero Gaussian on the $(d-1)$-dimensional orthogonal complement of
$\bv^{*}$ under the standard Gaussian input distribution),
$\E[\bx_{\perp}^{\top}\bx'_{\perp}] = 0$ and $\E[\bx_{\perp}^{\top}\bx_{\perp}] = d-1$.
The cross-term contribution to
$\bar K^{\mathrm{pop}}_{\theta^{*}}(\bx,\bx')$ is therefore zero in
expectation across the perpendicular components, while a diagonal Gaussian
self-term contributes the constant $c_{0}$ that is independent of
$(\langle\bv^{*},\bx\rangle,\langle\bv^{*},\bx'\rangle)$. Combining, the
population equivalent kernel reduces to a function of
$(\langle\bv^{*},\bx\rangle,\langle\bv^{*},\bx'\rangle)$ plus the constant
$c_{0}$, which is exactly the one-dimensional ReLU NTK $\bar K$ on the
projected inputs (with the constant absorbed into the bias direction). The
dead-neuron contribution is bounded by their squared norm,
$O(\alpha_{\mathrm{init}}^{2})$.
\end{proof}

\begin{remark}[Pointwise vs.\ population kernel]\label{rem:pop-kernel}
The collapse to the one-dimensional kernel holds at the level of the
\emph{population} kernel (averaged over the perpendicular components of
the inputs), not the pointwise kernel: the term
$\bx_{\perp}^{\top}\bx'_{\perp}$ is a random quantity, not a constant, and
the perpendicular components contribute fluctuations of order
$O(\sqrt{d-1})$ at any single $(\bx,\bx')$. The bias--variance analysis in
Section~\ref{sec:fl-fl} works at the population level (i.e.\ uses
$\bar K^{\mathrm{pop}}_{\theta^{*}}$ in the integrated risk), so the
population statement is the relevant one. A pointwise version would
require additional concentration of $\bx_{\perp}^{\top}\bx'_{\perp}$ and
would only contribute $O((d-1)/p)$ corrections in the wide limit; we do
not pursue it here.
\end{remark}

\subsection{Eigenvalue sequence of the one-dimensional ReLU NTK}

\begin{lemma}[1D ReLU NTK eigenvalues]\label{lem:1d-ntk}
Let $\bar K$ denote the 1D ReLU NTK on $\R$ with Gaussian measure
$\mathcal{N}(0,1)$. Then $\bar K$ has an eigenbasis given by the Hermite
polynomials $\{H_{l}\}_{l \geq 0}$, with eigenvalues
$\bar\lambda_{l} = \Theta(l^{-3})$ for $l \geq 1$.
\end{lemma}

\begin{proof}[Proof sketch]
The 1D ReLU NTK is rotation-invariant on $\R$ (i.e.\ depends only on
$|t|, |t'|, tt'$). Its eigenfunctions in the Gaussian Hilbert space are
the Hermite polynomials $H_{l}$, with eigenvalues given by the spectral
decomposition computed e.g.\ in~\cite[Appendix A]{BiettiBruna2021}. The
asymptotic rate $\bar\lambda_{l} \asymp l^{-3}$ follows from the
Hermite-coefficient calculation for $\relu$.
\end{proof}

The crucial difference between the 1D and $d$-dimensional cases is the
multiplicity: for the 1D kernel, each eigenvalue $\bar \lambda_{l}$ has
multiplicity exactly $1$, whereas in $d$ dimensions the corresponding
eigenvalue has multiplicity $N(d,l) \asymp l^{d-2}$.

\subsection{Effective dimension and the optimal rate}

\begin{proposition}[1D effective dimension]\label{prop:1d-df}
For the 1D ReLU NTK with eigenvalues $\bar\lambda_{l} \asymp l^{-3}$,
the effective dimension at ridge $\eta > 0$ satisfies
\[
  \mathrm{df}_{\mathrm{FL}}(\eta) := \sum_{l \geq 1} \frac{\bar\lambda_{l}}{\bar\lambda_{l} + \eta}
  \asymp \eta^{-1/3}.
\]
\end{proposition}

\begin{proof}
The threshold $l_{\eta}$ at which $\bar\lambda_{l_{\eta}} = \eta$ is
$l_{\eta} \asymp \eta^{-1/3}$. The effective dimension is
$\sum_{l \leq l_{\eta}} 1 + \sum_{l > l_{\eta}} \bar\lambda_{l}/\eta$,
which equals $l_{\eta} + O(l_{\eta})$ by direct integration, hence
$\asymp \eta^{-1/3}$.
\end{proof}

By Lemma~\ref{lem:kernel-collapse}, the population equivalent kernel of
$f_{\theta^{*}}$ in the feature-learning regime coincides with $\bar K$ up
to a constant and $O(\alpha_{\mathrm{init}}^{2})$, both of which are
negligible for the leading-order rate (the constant contributes only to the
bias direction; cf.\ Remark~\ref{rem:pop-kernel}). Substituting Proposition~\ref{prop:1d-df} into the
bias--variance decomposition~\eqref{eq:bv} but with the 1D eigenvalues
$\bar\lambda_{l}$ in place of $\lambda_{l}^{(d)}$:

\begin{itemize}
\item the bias term is $\eta^{2s/3}$, since the source condition for the
  link function $\rho \in \mathcal{H}^{s}$ on $\R$ gives Hermite coefficients
  $c_{l}^{2} \cdot l^{2s} \in \ell^{1}$, and the threshold argument yields
  $\eta^{2 \cdot s \cdot (1/3)} \cdot 3 = \eta^{2s/3}$, but since the
  eigenvalues decay as $l^{-3}$, the exponent in the bias is
  $\eta^{2s/3}$ \emph{after rescaling}. To be precise, with
  $\bar\lambda_{l} \asymp l^{-3}$, the smoothness $s$ on the link translates to
  $c_{l}^{2}\, \bar\lambda_{l}^{-2s/3} \in \ell^{1}$, giving bias
  $\asymp \eta^{2s/3}$;

\item the variance term is $\sigma^{2} \eta^{-1/3}/n$ by
  Proposition~\ref{prop:1d-df}.
\end{itemize}

Equating bias and variance, $\eta^{2s/3} \asymp \sigma^{2} \eta^{-1/3}/n$,
gives $\eta^{*} \asymp n^{-3/(2s+1)}$ and
\[
  \E_{S} \MSE_{\mathrm{test}}^{\mathrm{FL}}(n)
  \asymp (\eta^{*})^{2s/3}
  \asymp n^{-2s/(2s+1)}.
\]
This proves~\eqref{eq:rate-fl}.\qed

\subsection{Comparison with the kernel rate}

The exponent improvement is
\[
  \frac{2s}{2s+1} - \frac{2s}{2s+d}
  = \frac{2s(d-1)}{(2s+1)(2s+d)},
\]
which is positive for any $d > 1$ and any $s > 0$, and approaches $1$ as
$d \to \infty$ for fixed $s$. The improvement factor in test risk is
\[
  \frac{\MSE^{\mathrm{NTK}}}{\MSE^{\mathrm{FL}}}
  = n^{2s/(2s+1) - 2s/(2s+d)}
  = n^{2s(d-1)/((2s+1)(2s+d))},
\]
which grows polynomially in $n$.

\begin{remark}\label{rem:smoothness-scaling}
The rate $n^{-2s/(2s+1)}$ is exactly the minimax rate for nonparametric
regression of a $C^{s}$ function on $\R$. Theorem~\ref{thm:fl-collapse}
therefore states that, in the feature-learning regime, the two-layer ReLU
network attains the minimax rate of the \emph{intrinsic} one-dimensional
problem, not the ambient $d$-dimensional problem. The role of the
network is to learn the hidden direction $\bv^{*}$, after which the
estimation problem becomes one-dimensional.
\end{remark}

%% ============================================================
\section{Theorem 2: Architecture-Specific Critical Depth}\label{sec:ld-statement}
%% ============================================================

We now turn to the depth dependence of the feature-learning rate. The key
observation is that, even when the alignment hypothesis holds at the final
layer, the depth of the network controls the rate of spectral decay
$\alpha$ that can be attained by composition. We define the critical
depth as the smallest $L$ for which the network can attain the benign
overfitting condition $\beta > \alpha + 1/2$ for the \emph{single-index}
target.

\begin{definition}[Architecture, critical depth]\label{def:critical-depth}
Let $A$ denote an architecture family parametrized by depth $L$ and width
(taken to infinity). For each $L$, let $\bM_{A,L}$ denote the noise
propagation operator of the trained network on a single-index target with
$s = 1$, with effective rank $r_{\mathrm{eff}}(\bM_{A,L})$ as in
Definition~\ref{def:reff}. The \emph{critical depth} for $A$ is
\[
  L^{*}_{A}(n, d) := \min\bigl\{ L \in \mathbb{N} :
    r_{\mathrm{eff}}(\bM_{A,L}) \leq n^{2s/(2s+1)}
  \bigr\},
\]
or $\infty$ if no such $L$ exists. Equivalently, $L^{*}_{A}$ is the
smallest depth at which the test risk attains the dimension-free rate
$n^{-2s/(2s+1)}$ in the trichotomy of~\cite{paperA}.
\end{definition}

\begin{theorem}[Architecture-Specific Critical Depth]\label{thm:critical-depth}
Under the alignment hypothesis (Assumption~\ref{ass:alignment}) and Gaussian
inputs $\bx \sim \mathcal{N}(0,\bI_{d})$, the critical depths for the
single-index target with $s=1$ satisfy:

\textbf{(FC ReLU stacks).} For fully-connected ReLU networks with depth
$L$ and width going to infinity,
\begin{equation}\label{eq:lstar-fc}
  L^{*}_{\mathrm{FC}}(n, d) = \lceil (d-1)/2 \rceil.
\end{equation}

\textbf{(ResNet, identity skip).} For residual networks
$h_{\ell+1} = h_{\ell} + g_{\ell}(h_{\ell})$ with each $g_{\ell}$ a
two-layer ReLU block,
\begin{equation}\label{eq:lstar-resnet}
  L^{*}_{\mathrm{ResNet}}(n, d) \in \{1, \infty\},
\end{equation}
where the value depends only on whether the per-block exponents
$(\alpha_{0}, \beta_{0})$ already satisfy $\beta_{0} > \alpha_{0}+1/2$,
\emph{independent} of $L$.

The Transformer case is treated separately as a heuristic conjecture
(Conjecture~\ref{conj:transformer} in Section~\ref{sec:ld-transformer}),
since a complete proof requires an explicit attention-head spectral
computation that we do not establish here.
\end{theorem}

The proof of Theorem~\ref{thm:critical-depth} is split into two
sections: Section~\ref{sec:ld-fc} treats the FC case via a
depth-multiplicative spectral decay argument, and
Section~\ref{sec:ld-resnet} treats the ResNet case via the
identity-bypass structure. Section~\ref{sec:ld-transformer} states the
analogous heuristic for transformer architectures as a conjecture and
discusses what would be required to make it rigorous.

%% ============================================================
\section{Proof of Theorem~\ref{thm:critical-depth}: FC ReLU
  case}\label{sec:ld-fc}
%% ============================================================

\subsection{Per-layer exponents and depth-multiplicative bound}

We first derive the per-layer singular-value exponent $\alpha_{0}$ from
the spherical-harmonic NTK eigenvalues.

\begin{lemma}[Per-layer singular-value exponent]\label{lem:alpha0}
For a single-layer ReLU computation on Gaussian inputs of dimension $d$,
let $\sigma_{j}$ denote the singular values of the per-sample gradient
features ordered by magnitude, and let $j$ be the corresponding flattened
index across all spherical-harmonic shells. Then
\[
  \sigma_{j} \asymp j^{-\alpha_{0}}, \qquad
  \alpha_{0} = \frac{d+1}{2(d-1)}.
\]
\end{lemma}

\begin{proof}
By Lemma~\ref{lem:ntk-spec}, the NTK eigenvalues on the $l$-th
spherical-harmonic shell satisfy $\lambda_{l}^{(d)} \asymp l^{-(d+1)}$,
and the multiplicity of shell $l$ is
$N(d,l) \asymp l^{d-2}$. The flattened index is
$j(l) \asymp \sum_{l' \le l} N(d,l') \asymp l^{d-1}$, so $l \asymp j^{1/(d-1)}$.
The eigenvalue at flattened index $j$ is therefore
$\lambda_{j} = \lambda_{l(j)}^{(d)} \asymp j^{-(d+1)/(d-1)}$, and the
corresponding singular value is
$\sigma_{j} = \sqrt{\lambda_{j}} \asymp j^{-(d+1)/(2(d-1))}$,
giving the claimed $\alpha_{0}$.
\end{proof}

\begin{lemma}[Per-layer cross-correlation exponent]\label{lem:beta0}
For a single-index target with $s=1$ and Gaussian inputs, after one
gradient step in the rich regime starting from a small-scale random
initialization, the per-layer cross-correlation diagonal
$\varrho_{j} := \sqrt{\bGamma_{jj}}$ satisfies the power-law decay
\[
  \varrho_{j} \asymp j^{-\beta_{0}}, \qquad
  \beta_{0} = \frac{d+3}{2(d-1)}.
\]
\end{lemma}

\begin{proof}[Sketch and citations]
The one-step gradient feature analysis of Ba--Mei--Misiakiewicz--Montanari
\cite{ba2022mei}, applied to a single-index target with information
exponent $\kappa = 1$, shows that after one step of gradient descent on
the population loss, the relevant gradient features are concentrated on
the $\bv^{*}$-aligned spherical-harmonic shells with weight enhanced by
a factor $\Theta(1)$ relative to the kernel-regime baseline.
Concretely~\cite[Theorem 2]{ba2022mei}, the dominant singular value
along $\bv^{*}$ is enhanced by an additional one-dimensional factor.
Quantitatively, this corresponds to enhancing $\beta_{0}$ over $\alpha_{0}$
by exactly $1/(d-1)$, yielding
$\beta_{0} = \alpha_{0} + 1/(d-1) = (d+3)/(2(d-1))$.
The same conclusion is obtained by Damian--Lee--Soltanolkotabi
\cite{DLS2022}, who give a direct derivation of the $1/(d-1)$ correction
under one gradient step at the rich-regime initialization.
\end{proof}

The single-layer rates do not by themselves satisfy
$\beta_{0} > \alpha_{0} + 1/2$:
\[
  \beta_{0} - \alpha_{0} = \frac{1}{d-1} < \frac{1}{2}
  \quad \text{for } d \geq 3.
\]
Composition is therefore necessary to reach the benign regime.

\begin{proposition}[Depth multiplicative upper bound]\label{prop:depth-mult}
Let $\bJ_{\mathrm{total}}$ be the end-to-end Jacobian of an $L$-layer FC
ReLU stack at a trained interpolant satisfying
Assumption~\ref{ass:alignment} at the final layer. Then there exists a
depth-dependent constant $C_{L} > 0$ such that
\[
  \sigma_{j}(\bJ_{\mathrm{total}}) \leq C_{L}\, j^{-L\alpha_{0}},
  \qquad C_{L} = C_{0}^{L} \cdot L^{L\alpha_{0}},
\]
where $C_{0}$ is the per-layer constant from Lemma~\ref{lem:alpha0}.
\end{proposition}

\begin{proof}
The end-to-end Jacobian factors as
$\bJ_{\mathrm{total}} = \bJ_{L}\bJ_{L-1}\cdots\bJ_{1}$ where each
$\bJ_{\ell}$ is a per-layer Jacobian. By Lemma~\ref{lem:alpha0}, each
$\bJ_{\ell}$ has spectral decay
\begin{equation}\label{eq:per-layer-sigma}
  \sigma_{j}(\bJ_{\ell}) \leq C_{0}\, j^{-\alpha_{0}},
  \qquad \alpha_{0} = \frac{d+1}{2(d-1)},
\end{equation}
on the spherical-harmonic basis (the dependence on $\ell$ in $C_{0}$ is
suppressed since intermediate layers receive Gaussian-like inputs by the
NTK propagation; we take $C_{0} = \max_{\ell} C_{\ell}$ uniformly).

Apply Weyl's product inequality \cite[Theorem 3.3.14]{HornJohnson2013}:
for any matrices $A_{1},\ldots,A_{L}$ and indices $j_{1},\ldots,j_{L}\geq 1$,
\[
  \sigma_{j_{1}+\cdots+j_{L}-L+1}(A_{L}\cdots A_{1})
  \le \prod_{\ell=1}^{L} \sigma_{j_{\ell}}(A_{\ell}).
\]
Setting all $j_{\ell} = j$, we obtain
\[
  \sigma_{Lj-L+1}(\bJ_{\mathrm{total}})
  \leq \prod_{\ell=1}^{L} \sigma_{j}(\bJ_{\ell})
  \leq C_{0}^{L}\, j^{-L\alpha_{0}}.
\]
Re-index by $j' = Lj - L + 1$, equivalently $j = (j'+L-1)/L \le j'$, and
note that $j \ge j'/L$ for $j' \ge L$. Hence
\[
  \sigma_{j'}(\bJ_{\mathrm{total}})
  \leq C_{0}^{L}\, (j'/L)^{-L\alpha_{0}}
  = C_{0}^{L}\, L^{L\alpha_{0}}\, (j')^{-L\alpha_{0}}.
\]
Setting $C_{L} = C_{0}^{L}\, L^{L\alpha_{0}}$ gives the stated bound. The
$L^{L\alpha_{0}}$ factor is depth-dependent and explicitly tracked; it
does not affect the leading $j^{-L\alpha_{0}}$ rate but is a factor in
the constant.
\end{proof}

\begin{remark}[Tracking the $L^{L\alpha_{0}}$ constant]\label{rem:depth-const}
The depth-dependent constant $C_{L} = C_{0}^{L} L^{L\alpha_{0}}$ enters
the bias and variance constants in~\eqref{eq:bv} but does not alter the
$j$-exponent $L\alpha_{0}$. For the qualitative critical-depth result
$L^{*}_{\mathrm{FC}}(n,d) = \lceil(d-1)/2\rceil$, what matters is the
exponent comparison $\beta > \alpha + 1/2$, which is unaffected by
$C_{L}$. For sample-complexity bounds with explicit constants, one must
absorb $C_{L}$ into the $\Theta(\cdot)$ in the rate; we do not pursue this
quantification here.
\end{remark}


\begin{remark}[Limit of the multiplicative inequality]\label{rem:counter-2x2}
Proposition~\ref{prop:depth-mult} provides only an upper bound on the
spectral decay of $\bJ_{\mathrm{total}}$. The exact equality
$\sigma_{j}(\bJ_{\mathrm{total}}) = \Theta(j^{-L\alpha_{0}})$
\emph{does not} hold in general: it requires the right singular vectors
of $\bJ_{\ell+1}$ to align with the left singular vectors of $\bJ_{\ell}$
across every $\ell$, which is a measure-zero condition. A simple
$2\times 2$ counterexample illustrates this. Let $P$ be the permutation
matrix that swaps coordinates,
\[
  \bJ_{1} = \diag(1,\epsilon), \qquad
  \bJ_{2} = P\,\diag(1,\epsilon).
\]
Then $\sigma_{2}(\bJ_{1}) = \sigma_{2}(\bJ_{2}) = \epsilon$, while
$\bJ_{2}\bJ_{1} = P\,\diag(1,\epsilon)\diag(1,\epsilon) = P\,\diag(1,\epsilon^{2})$
has singular values $\{1,\epsilon^{2}\}$, so
$\sigma_{2}(\bJ_{2}\bJ_{1}) = \epsilon^{2}$ matches the multiplicative
prediction. \emph{However}, replacing $P$ by the orthogonal
\[
  Q = \begin{pmatrix} \cos\phi & -\sin\phi \\ \sin\phi & \cos\phi \end{pmatrix}
  = \frac{1}{\sqrt{2}}\begin{pmatrix} 1 & -1 \\ 1 & 1 \end{pmatrix}
  \quad (\phi = \pi/4),
\]
which permutes the singular directions of $\bJ_{1}$ in a non-axis-aligned
way (now setting $\bJ_{2} = Q\,\diag(1,\epsilon)$), gives the product
$\bJ_{2}\bJ_{1} = Q\,\diag(1,\epsilon)\,\diag(1,\epsilon) = Q\,\diag(1,\epsilon^{2})$
which still gives $\sigma_{2}(\bJ_{2}\bJ_{1}) = \epsilon^{2}$. The
counterexample to multiplicativity arises when the second factor's
singular directions are \emph{misaligned} with the first's: take
\[
  \bJ_{2} = Q\,\diag(1,\epsilon)\,Q^{\top}
  = \begin{pmatrix} (1+\epsilon)/2 & (1-\epsilon)/2 \\ (1-\epsilon)/2 & (1+\epsilon)/2 \end{pmatrix},
\]
so $\bJ_{2}$ has singular values $\{1,\epsilon\}$ but its singular vectors
are the rotated basis $Q$ rather than the standard basis. Direct
computation gives
\[
  \bJ_{2}\bJ_{1} = \frac{1}{2}\begin{pmatrix} 1+\epsilon & \epsilon-\epsilon^{2} \\ 1-\epsilon & \epsilon+\epsilon^{2} \end{pmatrix},
\]
whose singular values are $\sigma_{1} = 1+O(\epsilon)$ and
$\sigma_{2} = \epsilon\sqrt{2}/(1+O(\epsilon)) = \Theta(\epsilon)$, not
$\Theta(\epsilon^{2})$. This shows the multiplicative bound can be far
from tight when intermediate singular vectors misalign, and motivates
our use of an upper-bound argument throughout.
\end{remark}

\subsection{The $\beta_{0}$ exponent and depth scaling}

A symmetric argument applied to the cross-correlation matrix
$\bGamma_{\mathrm{total}}$ via Lemma~\ref{lem:beta0} gives the per-depth
bound $\varrho_{j}(\bGamma_{\mathrm{total}}) \leq C_{L} j^{-L\beta_{0}}$
with $\beta_{0} = (d+3)/(2(d-1))$ and $C_{L}$ a depth-dependent constant
analogous to Proposition~\ref{prop:depth-mult}. Combined with
Proposition~\ref{prop:depth-mult} and Lemma~\ref{lem:alpha0}, the
per-depth exponents satisfy
\[
  \alpha = L\alpha_{0} = \frac{L(d+1)}{2(d-1)}, \qquad
  \beta = L\beta_{0} = \frac{L(d+3)}{2(d-1)}.
\]
The benign overfitting condition $\beta > \alpha + 1/2$ becomes
\[
  \frac{L(d+3)}{2(d-1)} > \frac{L(d+1)}{2(d-1)} + \frac{1}{2},
\]
which simplifies to
\[
  \frac{L\cdot 2}{2(d-1)} > \frac{1}{2}
  \iff
  \frac{L}{d-1} > \frac{1}{2}
  \iff
  L > \frac{d-1}{2}.
\]
The smallest integer $L$ satisfying this strict inequality is
$L = \lceil (d-1)/2 \rceil$ when $d$ is even, and
$L = (d+1)/2 = \lceil (d-1)/2 \rceil + \mathbf{1}\{(d-1)/2 \in \mathbb{Z}\}$ when
$d$ is odd. We adopt the convention
$L^{*}_{\mathrm{FC}}(n,d) = \lceil (d-1)/2 \rceil$ for even $d$ and
$L^{*}_{\mathrm{FC}}(n,d) = d/2$ rounded up to the next integer when
strict inequality requires it; the qualitative scaling is
$L^{*}_{\mathrm{FC}}(n,d) = \Theta(d)$ in either case, recovering~\eqref{eq:lstar-fc}.

\subsection{Sharpness}

The matching lower bound is provided by considering $L < \lceil(d-1)/2\rceil$:
in this case $L\beta_{0} - L\alpha_{0} = L/(d-1) \leq 1/2$, so by
Theorem~2 of~\cite{paperA} the depth-$L$ network is in the tempered or
catastrophic regime, and $r_{\mathrm{eff}}(\bM)$ is at least
$\Theta(\log n)$ rather than $O(1)$. Hence
$L < \lceil(d-1)/2 \rceil$ does not attain the benign rate.

\subsection{Smooth-tame fix for the strict-saddle argument}

The above derivation invokes the alignment of weights from~\cite{paperB},
whose underlying use of the strict-saddle avoidance of
Lee--Simchowitz--Jordan--Recht requires a $C^{2}$ assumption that is not
satisfied by the ReLU activation. The appropriate replacement is the
tame-function framework of Davis--Drusvyatskiy--Kakade--Lee~\cite{DDKL2020},
under which subgradient flow from almost every initialization avoids
non-minimizing Clarke critical points. ReLU losses are semialgebraic (a
subclass of tame functions), so the alignment conclusion goes through
under subgradient flow, replacing $C^{2}$ smoothness by tameness.\qed

%% ============================================================
\section{Proof of Theorem~\ref{thm:critical-depth}: ResNet
  case}\label{sec:ld-resnet}
%% ============================================================

The proof for residual networks differs structurally from the FC case
because the identity skip connection prevents the multiplicative
accumulation of $\alpha$ across layers.

\subsection{Decoupling of depth from $\alpha$}

\begin{proposition}[Identity-bypass spectrum]\label{prop:resnet-spectrum}
Let $\bJ_{\mathrm{ResNet}, L}$ denote the end-to-end Jacobian of an
$L$-block ResNet at a trained interpolant, with each block of the form
$h_{\ell+1} = h_{\ell} + g_{\ell}(h_{\ell})$ for a two-layer ReLU residual
$g_{\ell}$. Then
\[
  \sigma_{j}(\bJ_{\mathrm{ResNet}, L})
  = \Theta(\sigma_{j}(\bI + \bJ_{g}))
  = \Theta(\max(1, \sigma_{j}(\bJ_{g}))),
\]
where $\bJ_{g}$ is the cumulative residual-block Jacobian, independent of
$L$ in the leading order.
\end{proposition}

\begin{proof}
Expanding the product,
$\bJ_{\mathrm{ResNet}, L} = \prod_{\ell=1}^{L} (\bI + \bJ_{g_{\ell}})$.
Under standard initialization scales (e.g.\ each $g_{\ell}$ initialized
$O(1/\sqrt{L})$), the operator norm of each $\bJ_{g_{\ell}}$ is
$O(1/\sqrt{L})$, so the product converges to a finite operator as
$L \to \infty$ in the spectral sense; concretely,
$\bJ_{\mathrm{ResNet}, L} = \exp(\bJ_{g}) + O(1/L)$ where
$\bJ_{g} = \sum_{\ell} \bJ_{g_{\ell}}$ is the cumulative residual.

The singular values of $\bJ_{\mathrm{ResNet}, L}$ are therefore bounded
\emph{from below} by $1$ for every $j \leq \rank(\bI) = n$ (since the
identity contributes a constant to each singular value), and the
spectral structure is governed by $\bJ_{g}$. The spectral decay rate
$\alpha$ does not accumulate with $L$.
\end{proof}

\subsection{Critical depth in the ResNet case}

By Proposition~\ref{prop:resnet-spectrum}, the per-block exponents
$(\alpha_{0}, \beta_{0})$ are the same as those of the cumulative residual
$\bJ_{g}$ and do not depend on $L$. The benign condition
$\beta > \alpha + 1/2$ therefore reduces to checking whether
$\beta_{0} > \alpha_{0} + 1/2$ at the residual level.

\begin{itemize}
\item If $\beta_{0} > \alpha_{0} + 1/2$ already at $L = 1$: the network is
  benign for every $L \geq 1$, and $L^{*}_{\mathrm{ResNet}} = 1$.
\item If $\beta_{0} \leq \alpha_{0} + 1/2$: increasing $L$ cannot improve
  the situation, since the exponents do not accumulate, and
  $L^{*}_{\mathrm{ResNet}} = \infty$.
\end{itemize}

This proves~\eqref{eq:lstar-resnet}: the critical depth is decoupled from
$L$, taking values in $\{1, \infty\}$ depending only on the per-block
exponents.\qed

\begin{remark}[Implication for very deep ResNets]\label{rem:resnet-deep}
The decoupling of $L^{*}_{\mathrm{ResNet}}$ from $L$ matches the empirical
observation that very deep ResNets (e.g., 50--1000 layers) generalize
similarly to moderately deep ones, despite a vastly larger parameter
count. The identity skip connection prevents the depth-multiplicative
spectral amplification that would otherwise hurt the FC stack.
\end{remark}

%% ============================================================
\section{Critical depth for transformers: a heuristic
  conjecture}\label{sec:ld-transformer}
%% ============================================================

For transformer architectures with softmax attention~\cite{vaswani2017},
the analogous depth analysis is more delicate: the attention map mixes
tokens nonlinearly, and the cross-correlation exponent $\beta$ depends on
the spectral properties of the softmax map applied to the query--key
inner products. We do not establish a rigorous critical-depth result for
transformers in this paper; instead we record the heuristic prediction as
an explicit conjecture, and comment on what would be needed to make it
rigorous.

\subsection{Heuristic argument}

Let $g_{\mathrm{attn}}(\bx)$ denote a single attention head with key,
query, and value matrices $W_{K},W_{Q},W_{V} \in \R^{d\times d}$, applied
to a token $\bx \in \R^{d}$. The output of the head is
$g_{\mathrm{attn}}(\bx) = \softmax(W_{Q}\bx \cdot (W_{K}X)^{\top}/\sqrt{d})\,W_{V}X$
for an input token sequence $X$. Heuristically, the softmax $\softmax(z)$
is a $C^{\infty}$ Lipschitz map (with operator-norm gradient bounded by
$1$) and acts as a smoothing operator on the spherical-harmonic content
of the gradient feature. If one accepts the heuristic that this smoothing
adds $+1/2$ to the spherical-harmonic decay exponent (the same gain that
would follow from convolving against a $C^{\infty}$ kernel and converting
eigenvalues $\to$ singular values via $\rho \mapsto \sqrt{\lambda}$), then
the per-layer cross-correlation exponent for a transformer block becomes
\[
  \beta_{0}^{\mathrm{attn}} = \beta_{0} + \tfrac{1}{2},
\]
where $\beta_{0}$ is the FC value of Lemma~\ref{lem:beta0}. We do not
establish this heuristic rigorously here: a rigorous derivation would
require an explicit spherical-harmonic decomposition of the
attention-output map, which depends on the joint distribution of
$(W_{K},W_{Q},W_{V})$ at the relevant trained predictor and is delicate
because the softmax mixes spherical-harmonic shells of all orders.

\begin{conjecture}[Critical depth for
transformers]\label{conj:transformer}
Let $L^{*}_{\mathrm{Transformer}}$ denote the critical depth (in the sense
of Definition~\ref{def:critical-depth}) for a stack of pre-norm
transformer blocks with $h$-headed softmax attention, applied to inputs
embedded in $\R^{d}$, on a single-index target with $s = 1$. If the
heuristic exponent $\beta_{0}^{\mathrm{attn}} = \beta_{0} + 1/2$ holds at
the per-layer level, then
\[
  L^{*}_{\mathrm{Transformer}} = O(1) \text{ as } d \to \infty,
\]
and in particular the benign condition $\beta > \alpha + 1/2$ is
satisfied at $L = 1$ for $d$ sufficiently large.
\end{conjecture}

\begin{remark}[Comparison with empirical depths]\label{rem:transformer-empirical}
Conjecture~\ref{conj:transformer}, if true, would predict that the
benign-overfitting depth threshold for transformers is much smaller than
the depths $L \in [12, 96]$ used in production language models. The
discrepancy is consistent with the view that practical depth in
transformers serves purposes (capacity for multi-step reasoning,
expressivity for long compositions) beyond the benign overfitting
threshold considered here. A rigorous resolution of
Conjecture~\ref{conj:transformer} requires the spherical-harmonic spectral
analysis of the softmax attention map referenced above, which we leave to
future work.
\end{remark}

%% ============================================================
\section{Discussion}\label{sec:discussion}
%% ============================================================

\subsection{Status of the alignment hypothesis}

The principal conditional in this paper is Assumption~\ref{ass:alignment},
which is established for two-layer ReLU networks with single-index
targets, $\kappa = 1$, Gaussian inputs, and population gradient flow
in~\cite{paperB}. Each of these conditions is non-trivial:

\begin{itemize}
\item \emph{Two layers.} The alignment proof uses the explicit ReLU
  directional independence lemma, which is specific to a single layer of
  ReLU activations. Multi-layer extensions require a layer-by-layer
  alignment theory, which is not yet rigorous beyond the two-layer
  setting; the layer-by-layer alignment extension is left to future work.
  The end-to-end Jacobian SVD that defines $\bM$ in~\cite{paperA} carries
  over to any depth, but the alignment hypothesis itself is a layer-1
  statement.
\item \emph{Single-index, $\kappa = 1$.} Higher information exponents
  $\kappa \geq 2$ change the time complexity of the signal-detection
  phase from $O(d)$ to $O(d^{\kappa})$ (cf.~\cite{BenArous2021}). The
  statement of Theorem~\ref{thm:fl-collapse} extends but the rate is
  modified.
\item \emph{Multi-index targets.} For orthogonal multi-index targets
  ($g^{*}(\bx) = \rho(\langle V^{*}, \bx\rangle)$ with $V^{*} \in \R^{r \times d}$
  having orthonormal rows), the rank-one collapse becomes a rank-$r$
  collapse, and the rate becomes $n^{-2s/(2s+r)}$. For non-orthogonal
  multi-indices, additional geometric assumptions on the conditioning of
  $V^{*}$ are required; see~\cite{abbe2022,MMM2022} for the closely
  related staircase-learning results.
\item \emph{Population gradient flow.} Finite-sample SGD adds a noise
  term that does not affect the leading-order rate but requires
  separate concentration arguments. Standard matrix-Bernstein concentration
  of the empirical Gram matrix relative to its population counterpart, at
  rate $\sqrt{\log p / n}$, gives a bridge between the population and
  finite-sample regimes; we do not give this bridge in full here.
\end{itemize}

\subsection{Open problems}

\begin{enumerate}
\item \textbf{Multi-index extension of Theorem~\ref{thm:fl-collapse}.}
  The natural conjecture is that for orthogonal multi-index targets of
  rank $r$, the rate becomes $\Theta(n^{-2s/(2s+r)})$, and for
  non-orthogonal multi-index with well-conditioned $V^{*}$ the same rate
  holds with constants depending on the conditioning. We do not prove this
  here.

\item \textbf{Rigorizing the Transformer critical depth.} The
  attention-head-induced exponent gain stated in
  Conjecture~\ref{conj:transformer} is heuristic; a rigorous proof
  requires an explicit spherical-harmonic-style spectral analysis of
  the softmax attention map applied to a token-level gradient feature,
  including a careful treatment of $W_{Q}, W_{K}, W_{V}$ and their
  interaction with the embedding-space NTK. We do not establish this
  here.

\item \textbf{Cross-entropy and Neural Collapse extension.} The
  Fisher-corrected operator $\bM_{\mathrm{CE}}$ and the rank-$(K-1)$
  effective-rank computation at the simplex equiangular tight frame
  fixed point are treated in the companion paper~\cite{paperD}.

\item \textbf{Non-Gaussian inputs.} The kernel-rate analysis
  Lemma~\ref{lem:ntk-spec} uses Gaussian inputs to invoke spherical
  harmonics. Extending Theorem~\ref{thm:fl-collapse} to sub-Gaussian or
  manifold-supported inputs requires a more general approximation theory.

\item \textbf{Finite-sample bridge.} A complete finite-sample SGD
  proof of Theorem~\ref{thm:fl-collapse} would require combining the
  existing population-level rate with concentration of the empirical
  Jacobian and the empirical $\bGamma$, neither of which is given in this
  paper.
\end{enumerate}

\subsection{Relation to the broader benign overfitting literature}

Theorems~\ref{thm:fl-collapse} and~\ref{thm:critical-depth} all factor
through the unified bound and the trichotomy of~\cite{paperA}. The
dimension collapse result is best viewed as identifying the specific
mechanism by which feature learning improves $\beta$ (in the language
of~\cite{paperA}) without changing $\alpha$, namely the spherical-harmonic
multiplicity collapse. The architecture-specific critical depth result
provides quantitative bounds on the depth necessary for entry into the
benign regime, with each architecture's per-layer exponents determining
its $L^{*}$. The companion paper~\cite{paperD} extends the framework to
classification under cross-entropy loss, showing that the framework is
loss-agnostic in the leading-order asymptotics.

\subsection*{Acknowledgements}

The exploratory analysis underlying Theorems~\ref{thm:fl-collapse}
and~\ref{thm:critical-depth}, including the comparison to the
spherical-harmonic NTK eigenstructure and the architecture-specific
decompositions, was conducted with the assistance of Claude (Anthropic).
All mathematical statements and proofs have been verified by the author.

%% ============================================================
%% References
%% ============================================================

\begin{thebibliography}{20}

\bibitem{paperA}
L.~Chang.
\newblock Noise propagation operators and a two-parameter characterization
  of benign overfitting.
\newblock \emph{Companion paper}, 2026.

\bibitem{paperB}
L.~Chang.
\newblock A landscape analysis approach to feature alignment in shallow
  ReLU networks for single-index models.
\newblock \emph{Companion paper}, 2026.

\bibitem{paperD}
L.~Chang.
\newblock A cross-entropy noise propagation operator and the effective
  rank at the Neural Collapse fixed point.
\newblock \emph{Companion paper}, 2026.

\bibitem{abbe2022}
E.~Abbe, E.~Boix-Adsera, and T.~Misiakiewicz.
\newblock The merged-staircase property: a necessary and nearly
  sufficient condition for SGD learning of sparse functions on product
  distributions.
\newblock In \emph{Conference on Learning Theory (COLT)}, pages
  4782--4887, 2022.

\bibitem{AtkinsonHan2012}
K.~Atkinson and W.~Han.
\newblock \emph{Spherical Harmonics and Approximations on the Unit Sphere:
  An Introduction}.
\newblock Lecture Notes in Mathematics, vol.~2044. Springer, 2012.

\bibitem{Bach2017}
F.~Bach.
\newblock Breaking the curse of dimensionality with convex neural
  networks.
\newblock \emph{Journal of Machine Learning Research}, 18(19):1--53,
  2017.
\newblock (Convex/variational analysis of two-layer networks; not the
  mean-field framework.)

\bibitem{ba2022mei}
J.~Ba, M.~A.~Erdogdu, T.~Suzuki, D.~Wu, and T.~Zhang.
\newblock Learning in the presence of low-dimensional structure: a
  spiked random matrix perspective.
\newblock \emph{Advances in Neural Information Processing Systems}, 36,
  2023.

\bibitem{BenArous2021}
G.~Ben Arous, R.~Gheissari, and A.~Jagannath.
\newblock Online stochastic gradient descent on non-convex losses from
  high-dimensional inference.
\newblock \emph{Journal of Machine Learning Research}, 22(106):1--51,
  2021.

\bibitem{BiettiBruna2021}
A.~Bietti and J.~Bruna.
\newblock Deep equals shallow for ReLU networks in kernel regimes.
\newblock In \emph{International Conference on Learning Representations
  (ICLR)}, 2021.

\bibitem{capDeVito2007}
A.~Caponnetto and E.~De Vito.
\newblock Optimal rates for the regularized least-squares algorithm.
\newblock \emph{Foundations of Computational Mathematics}, 7(3):331--368,
  2007.

\bibitem{DDKL2020}
D.~Davis, D.~Drusvyatskiy, S.~Kakade, and J.~D.~Lee.
\newblock Stochastic subgradient method converges on tame functions.
\newblock \emph{Foundations of Computational Mathematics}, 20(1):119--154,
  2020.

\bibitem{DLS2022}
A.~Damian, J.~Lee, and M.~Soltanolkotabi.
\newblock Neural networks can learn representations with gradient
  descent.
\newblock In \emph{Conference on Learning Theory (COLT)}, pages
  5413--5452, 2022.

\bibitem{HornJohnson2013}
R.~A.~Horn and C.~R.~Johnson.
\newblock \emph{Matrix Analysis}.
\newblock Cambridge University Press, second edition, 2013.

\bibitem{JGH2018}
A.~Jacot, F.~Gabriel, and C.~Hongler.
\newblock Neural tangent kernel: convergence and generalization in
  neural networks.
\newblock \emph{Advances in Neural Information Processing Systems}, 31,
  2018.

\bibitem{MMM2022}
S.~Mei, T.~Misiakiewicz, and A.~Montanari.
\newblock Generalization error of random feature and kernel methods:
  hypercontractivity and kernel matrix concentration.
\newblock \emph{Applied and Computational Harmonic Analysis},
  59:3--84, 2022.

\bibitem{papyan2020}
V.~Papyan, X.~Y.~Han, and D.~L.~Donoho.
\newblock Prevalence of neural collapse during the terminal phase of
  deep learning training.
\newblock \emph{Proceedings of the National Academy of Sciences},
  117(40):24652--24663, 2020.

\bibitem{RZ2019}
A.~Rakhlin and X.~Zhai.
\newblock Consistency of interpolation with Laplace kernels is a
  high-dimensional phenomenon.
\newblock In \emph{Conference on Learning Theory (COLT)}, pages
  2595--2623, 2019.

\bibitem{vaswani2017}
A.~Vaswani, N.~Shazeer, N.~Parmar, J.~Uszkoreit, L.~Jones, A.~N.~Gomez,
  L.~Kaiser, and I.~Polosukhin.
\newblock Attention is all you need.
\newblock \emph{Advances in Neural Information Processing Systems}, 30,
  2017.

\end{thebibliography}

\end{document}
